
\input texsis
\paper

\overfullrule=0pt
\def\frac#1#2{#1\over #2}
%
\def\bmatrix#1{\left[ \matrix{#1} \right]}
\def\be{{\beta}}
\def\toone{{t+1}}

%\def\frac#1#2{#1\over #2}
%\magnification=1200
\overfullrule=0pt
%
 

\centerline{\bf Lucas II}


We let $s_t \in \{1,2\}$ be governed by a Markov chain with transition matrix $P$ where
$P_{ij} = {\rm Prob}(s_{t+1} = j| s_t  = i) $ and initial probability distribution $\pi_0$.
We assume that $\{s_t, x_t{t+1}, y_t\}_{t=0}^\infty$ is governed by  the Markov-switching linear state-space system
$$ \eqalign{ x_{t+1} & = A_{s_t} x_{t+1} + C_{s_t} w_{t+1} \cr
                    y_t  & = G_{s_t} x_t  \cr  } $$
where $x_0 $ is a random initial vector drawn from  density  $\phi_0$ and $w_{t+1} \sim {\cal N}(0,I)$ is an i.i.d.~sequence of random
vectors.                      
We   want to compute the conditional  mathematical expectations
 $$ z_t = E ( \sum_{t=0}^\infty \beta^t y_{t+j} ) \vert x_t, s_t  $$     
 and 
 $$ \tilde z_t =  E ( \sum_{t=0}^\infty \beta^t y_{t+j} ) | x_t . $$    
 
We compute $z_t$ first.        
We begin by guessing that
$$ z_t = H_{s_t} x_t  $$
and substitute this guess into
$$ z_t = y_t + \beta E z_{t+1} | x_t, s_t  $$
to get 
$$  H_i  = G_{i}   + \beta \sum_{j} H_j A_i P_{ij}. $$
which implies in particular that
$$ \eqalign{ H_1 & = G_1 + \beta [ H_1 A_1 P_{11} + H_2 A_1 P_{12} ] \cr
                   H_2 & = G_2 + \beta [ H_1 A_2 P_{21} + H_2 A_1 P_{22} ]. }$$                  
Stacking and solving, we find
$$ \bmatrix{H_1 & H_2 }  = \bmatrix{G_1 & G_2}             (I - \beta \bmatrix {A_1 P_{11} & A_2 P_{21} \cr
                                                   A_1 P_{12} & A_2 P_{22} }  )^{-1}. $$                   


We now move on to computing $ \tilde z_t =  E ( \sum_{t=0}^\infty \beta^t y_{t+j} ) | x_t . $    
We let $\pi(t)$ be a distribution over the finite Markov state $s_t$ and denote $\pi_i(t)$ be the probability
of state $i$ at $t$.  For now, we'll suppress the $t$ and just write $\pi_i$ with $t$ being understood, depending on the
context of our application.
Note that the 
 following  normal distributions  conditional on $x_t, s_t$ are multivariate normal:
$$ \eqalign{ x_{t+1} | x_t, s_t  & \sim {\cal N}( A_{s_t} x_t, C_{s_t}C_{s_t} ' ) \cr
                    y_t | x_t, s_t & \sim {\cal N}( G_{s_t} x_t, 0 )  \cr
                    z_t | x_t, s_t & \sim {\cal N}( H_{s_t} x_t, 0 ) }  $$
But distributions $x_{t+1}, y_t$, and $z_t$ conditional on $x_t$ are {\it mixtures of normals}, with $\pi_{i,t}$  the distribution over $s_t$ being the mixing probabilities over states $i$  at time $t$.              
Defining the matrix averages
$$ \eqalign{ \bar A_t & = \pi_{1,t} A_1 + \pi_{2,t} A_2   \cr
                    \bar G_t & =\pi_{1,t} G_1 + \pi_{2,t} G_2 \cr
                    \bar H_t & = \pi_{1,t} H_1 + \pi_{2,t} H_2 }$$ 
we can deduce three conditional means
$$ \eqalign{ E x_{t+1} | x_t & = \bar A_t x_t \cr
                    E y_t | x_t & = \bar G_t x_t      \cr
                    E z_t | x_t & = \bar H_t x_t }$$ 
The appearance of the time subscripts $t$ in the formulas for $\bar A_t, G_t, H_t$ are a consequence of  $\pi_{i,t}$ not necessarily
being a stationary distribution for Markov chain $P$.                
                    
                    
Now let's compute conditional covariances.  
Since
$$ x_{t+1} - \bar A x_t = \cases{ (A_1 - \bar A_t) x_t  + C_1 w_{t+1} & if $s_t =1$ \cr
                                                    (A_2 - \bar A_t) x_t  + C_2 w_{t+1} & if $s_t =2$ } $$
it follows that 
$$ E (x_{t+1} - \bar A_t x_t ) (x_{t+1} - \bar A _tx_t ) ' | x_t = \sum_i \pi_{i,t} [ (A_i - \bar A_t) x_t x_t' (A_i - \bar A_t)' + C_i C_i' ] $$
And since                  
$$ y_t = \cases{ G_1 x_t & if $s_t =1$ \cr
                           G_2 x_t & if $s_t =2 $ } $$                     
and 
$$ y_t - \bar G x_t = \cases{ (G_1 - \bar G_t) x_t & if $s_t =1$ \cr
                           (G_2  - \bar G_t) x_t & if $s_t =2 $ } $$                             
it follows that 
$$ E (y_t - \bar G_t x_t) (y_t - \bar G_t x_t)' | x_t = \sum_i \pi_{i,t} (G_i - \bar G_t) x_t x_t' (G_i - \bar G_t)'.  $$   
and that
$$ E (z_t - \bar H_t x_t) (z_t - \bar H_t x_t)' | x_t = \sum_i \pi_{i,t} (H_i - \bar H_t) x_t x_t' (H_i - \bar H_t)'.  $$                                                    
       
                
                    
                                              
Let 
$$\Sigma(x_t)  = \sum_i \pi_i [ (A_i - \bar A) x_t x_t' (A_i - \bar A)' + C_i C_i' ] $$
Then since $\{x_{t}\}$ is a stationary and ergodic process with stationary distribution $\pi(x)$, we have a mean-ergodic theorem that
tells us that as $ {T \rightarrow + \infty}$
$$ T^{-1} \sum_{t=0}^T \Sigma(x_t)  \rightarrow \int \Sigma(x) d \pi(x) \equiv \Sigma $$
A vector autoregression of the form
$$ x_{t+1} = \bar A x_t + v_{t+1}$$
has $E v_{t+1} v_{t+1}' = \Sigma$.

We can compute  means and covariances directly as follows.
Let
$E x_t = \mu_t$ and $\Sigma_t = E (x_t - \mu_t) (x_t  - \mu_t)'$.  
Note that
$$ x_{t+1} - \mu_{t+1} = \cases{ (A_1 - \mu_t) x_t  + C_1 w_{t+1} & if $s_t =1$ \cr
                                                    (A_2 - \mu_t) x_t  + C_2 w_{t+1} & if $s_t =2$ } $$
It follows that
$$ \mu_{t+1} = \bar A_t \mu_t$$
and
$$\Sigma_{t+1} = \pi_1 (A_1 - \bar A_t) \Sigma_t (A_1 - \bar A_t)' + \pi_2 (A_2 - \bar A_t) \Sigma_t (A_2 - \bar A_t)' + \bar C \bar C ' $$
and that iterations on these two recursions will converge, respectively, to 
a stationary mean vector $\mu$  and a stationary covariance matrix $\Sigma$ that
satisfy
$$ \mu  = \bar A \mu $$
and
$$\Sigma = \pi_1 (A_1 - \bar A) \Sigma (A_1 - \bar A)' + \pi_2 (A_2 - \bar A) \Sigma(A_2 - \bar A)' + \bar C \bar C ' .$$


\subsection{Example}

Let $p_t$ be the rate  of inflation and $m_t$ the rate of growth of the money supply, where the rate of inflation is the first difference of the logarithm of the price level
and the rate of growth of the money supply is the first difference of the logarithm of the money supply.  Assume
that 
$$ p_t = (1-\lambda) m_t + \lambda E_t p_{t+1},  \quad \lambda \in (0,1)  $$
so that
$$ p_t = (1-\lambda) E_t \sum_{j=0}^\infty \lambda^j m_{t+j} , $$
where $E_t (\cdot)$ is a mathematical expectation conditioned on time $t$ information, which will be either $x_t$ or $x_t, s_t$ where $x_t, s_t$ are to be defined now.
We assume  that $s_t \in {1,2}$ is governed by a two-state Markov chain $P$ with initial distribution $\pi_0$. 
As for $x_t$, it is a state that lets us represent the following stochastic process for $\{m_t\}_{t=0}^\infty$ in the state-space form XXXX:
$$ m_{t+1} = \rho_{1,s(t)} m_t + \rho_{2,s(t)} m_{t-1} + c_{s(t)}w_{t+1} $$
where $w_{t+1}$ is an i.i.d. standardized normal random variable. 

\noindent{\bf Quentin:} Please compute ``all of our objects'' for this example with two $s_t$ dependent   $\rho_1, \rho_2$ vectors, namely $1.2, -.3$ and $.95, 0$ with
a transition matrix $P = \bmatrix{ .5 & .5 \cr .5 & .5 } $ using as $\pi_0 =\pi_\infty$ the stationary distribution $\bmatrix{.5 & .5}'$







                                                    







\bye
